\chapter{Conclusion}

The transition from traditional Public-Key Encryption (PKE) and interactive Diffie-Hellman paradigms to the Key Encapsulation Mechanism (KEM) represents one of the most profound architectural shifts in modern applied cryptography. This paper has traced the complete trajectory of KEM, demonstrating how a simple syntactic refinement---delegating the encapsulation of a random secret to the asymmetric primitive---has streamlined security proofs, enabled generic composition, and paved the way for the post-quantum era.

Through our technical review, we highlighted the critical milestones that matured KEM into a robust standard. The KEM/DEM systematization and the Cramer-Shoup construction established IND-CCA security as the uncompromising baseline for modern network protocols. Subsequently, the modular Fujisaki-Okamoto (FO) transformation provided the essential bridge for post-quantum cryptography, allowing inherently noisy and CPA-secure lattice-based primitives to achieve CCA security safely and efficiently.

To contextualize these theoretical constructs, we developed and evaluated jkem, a standard-compliant Rust implementation of ML-KEM. Our empirical evaluation emphasizes that the security of modern KEMs heavily relies not only on mathematical hardness assumptions but also on meticulous software engineering. Implementing constant-time decapsulation and correctly managing the implicit rejection paths are as crucial as the underlying Module-LWE problem itself. While portable software implementations offer excellent educational and structural clarity, our benchmarks reaffirm the necessity for architecture-specific optimizations to meet the stringent latency requirements of global-scale protocols.

Finally, our analysis of HPKE and Hybrid TLS 1.3 deployments underscores a fundamental reality: a KEM is a powerful building block, but it is not a standalone secure channel. The overarching protocols must rigorously enforce context binding, identity authentication, and downgrade resistance. As the cryptographic ecosystem migrates toward post-quantum resilience, the true success of KEM will be measured by our ability to seamlessly integrate these encapsulation mechanisms into hybrid architectures, balancing conservative security margins with computational efficiency. Ultimately, the KEM abstraction has proven to be the right interface at the right time, providing the foundational stability needed to secure tomorrow's digital infrastructure against quantum threats.